\subsection{Quản lý Chapter (Chapter Management)}\label{subsec:srs_chapter_mgmt}
\subsubsection{UC-11: Tạo Chapter mới}\label{subsubsec:uc11_create_chapter}
\begin{table}[H]
\centering
\renewcommand{\arraystretch}{1.3}
\begin{tabularx}{\textwidth}{|>{\raggedright\arraybackslash\bfseries}p{4.5cm}|>{\raggedright\arraybackslash}X|}
\hline
\rowcolor{tableheader}
\multicolumn{2}{|c|}{\textbf{Đặc tả Use Case: Tạo Chapter mới}} \\ \hline
\textbf{Mã Use Case} & UC-11 \\ \hline
\textbf{Tên Use Case} & Tạo Chapter mới \\ \hline
\textbf{Actor chính} & Mangaka \\ \hline
\textbf{Mô tả} & Cho phép Mangaka tạo một Chapter mới thuộc về một Series cụ thể. \\ \hline
\textbf{Tiền điều kiện (Pre-conditions)} & Tài khoản Mangaka ở trạng thái Active và đã được cấp quyền quản lý nội dung của Series tương ứng. \\ \hline
\textbf{Luồng sự kiện chính (Basic Flow)} & 
\begin{minipage}[t]{\linewidth}
    \begin{enumerate}[leftmargin=*]
        \item Mangaka chọn Series cần thêm Chapter.
        \item Mangaka chọn chức năng 'Thêm Chapter mới'.
        \item Hệ thống hiển thị form nhập liệu.
        \item Mangaka nhập thông tin (Số thứ tự lớn hơn 0, Tiêu đề tối đa 100 ký tự, Số trang dự kiến).
        \item Mangaka nhấn 'Lưu'.
        \item Hệ thống xác thực dữ liệu, lưu thông tin và hiển thị thông báo thành công.
    \end{enumerate}
\end{minipage} \\ \hline
\textbf{Luồng rẽ nhánh/Ngoại lệ (Alternate/Exception)} & 
\begin{minipage}[t]{\linewidth}
    \begin{itemize}[leftmargin=*]
        \item \textbf{A1}: Thông tin không hợp lệ (trùng số thứ tự). Hệ thống báo lỗi chữ đỏ và tự động gợi ý số thứ tự Chapter hợp lệ kế tiếp.
    \end{itemize}
\end{minipage} \\ \hline
\textbf{Hậu điều kiện (Post-conditions)} & Một Chapter mới được tạo trong hệ thống với trạng thái 'Bản nháp' (Draft). \\ \hline
\end{tabularx}
\caption{Đặc tả UC-11: Tạo Chapter mới}
\label{tab:uc11}
\end{table}

\subsubsection{UC-12: Cập nhật thông tin Chapter}\label{subsubsec:uc12_update_chapter}
\begin{table}[H]
\centering
\renewcommand{\arraystretch}{1.3}
\begin{tabularx}{\textwidth}{|>{\raggedright\arraybackslash\bfseries}p{4.5cm}|>{\raggedright\arraybackslash}X|}
\hline
\rowcolor{tableheader}
\multicolumn{2}{|c|}{\textbf{Đặc tả Use Case: Cập nhật thông tin Chapter}} \\ \hline
\textbf{Mã Use Case} & UC-12 \\ \hline
\textbf{Tên Use Case} & Cập nhật thông tin Chapter \\ \hline
\textbf{Actor chính} & Mangaka \\ \hline
\textbf{Mô tả} & Cho phép Mangaka chỉnh sửa thông tin của Chapter đã tạo (tiêu đề, số trang dự kiến). \\ \hline
\textbf{Tiền điều kiện (Pre-conditions)} & Mangaka có quyền quản lý Series và Chapter phải đang ở trạng thái cho phép chỉnh sửa (chưa xuất bản). \\ \hline
\textbf{Luồng sự kiện chính (Basic Flow)} & 
\begin{minipage}[t]{\linewidth}
    \begin{enumerate}[leftmargin=*]
        \item Mangaka chọn Chapter cần chỉnh sửa.
        \item Hệ thống hiển thị thông tin hiện tại của Chapter.
        \item Mangaka thay đổi thông tin.
        \item Mangaka nhấn 'Cập nhật'.
        \item Hệ thống kiểm tra tính hợp lệ, lưu thay đổi và hiển thị thông báo thành công.
    \end{enumerate}
\end{minipage} \\ \hline
\textbf{Luồng rẽ nhánh/Ngoại lệ (Alternate/Exception)} & 
\begin{minipage}[t]{\linewidth}
    \begin{itemize}[leftmargin=*]
        \item \textbf{A1}: Chapter đã xuất bản. Hệ thống khóa (disable) các trường quan trọng (số thứ tự) và hiển thị thông báo 'Chỉ được sửa tiêu đề phụ'.
    \end{itemize}
\end{minipage} \\ \hline
\textbf{Hậu điều kiện (Post-conditions)} & Thông tin Chapter được cập nhật đồng bộ trong cơ sở dữ liệu. \\ \hline
\end{tabularx}
\caption{Đặc tả UC-12: Cập nhật thông tin Chapter}
\label{tab:uc12}
\end{table}

\subsubsection{UC-13: Theo dõi tiến độ Chapter}\label{subsubsec:uc13_track_chapter}
\begin{table}[H]
\centering
\renewcommand{\arraystretch}{1.3}
\begin{tabularx}{\textwidth}{|>{\raggedright\arraybackslash\bfseries}p{4.5cm}|>{\raggedright\arraybackslash}X|}
\hline
\rowcolor{tableheader}
\multicolumn{2}{|c|}{\textbf{Đặc tả Use Case: Theo dõi tiến độ Chapter}} \\ \hline
\textbf{Mã Use Case} & UC-13 \\ \hline
\textbf{Tên Use Case} & Theo dõi tiến độ Chapter \\ \hline
\textbf{Actor chính} & Mangaka, Tantou Editor \\ \hline
\textbf{Mô tả} & Xem tổng quan tiến độ hoàn thành các trang và công việc (Task) trong Chapter. \\ \hline
\textbf{Tiền điều kiện (Pre-conditions)} & Tài khoản ở trạng thái Active và có quyền xem báo cáo của Series tương ứng. \\ \hline
\textbf{Luồng sự kiện chính (Basic Flow)} & 
\begin{minipage}[t]{\linewidth}
    \begin{enumerate}[leftmargin=*]
        \item Người dùng chọn Chapter cần theo dõi.
        \item Người dùng chọn 'Tiến độ'.
        \item Hệ thống tính toán và hiển thị \% hoàn thành dựa trên số Task/Page đã chuyển sang trạng thái 'Completed'.
        \item Hiển thị danh sách các trang đang thực hiện, đang chờ duyệt bằng biểu đồ trực quan.
    \end{enumerate}
\end{minipage} \\ \hline
\textbf{Luồng rẽ nhánh/Ngoại lệ (Alternate/Exception)} & 
\begin{minipage}[t]{\linewidth}
    \begin{itemize}[leftmargin=*]
        \item \textbf{A1}: Chapter chưa có Page hay Task nào. Hệ thống hiển thị 0\% và gợi ý nút 'Tạo Task ngay' hoặc 'Tạo Page mới'.
    \end{itemize}
\end{minipage} \\ \hline
\textbf{Hậu điều kiện (Post-conditions)} & Người dùng nắm được tình trạng hiện tại của Chapter theo thời gian thực. \\ \hline
\end{tabularx}
\caption{Đặc tả UC-13: Theo dõi tiến độ Chapter}
\label{tab:uc13}
\end{table}

\subsection{Quản lý Page (Page Management)}\label{subsec:srs_page_mgmt}
\subsubsection{UC-14: Tạo Page}\label{subsubsec:uc14_create_page}
\begin{table}[H]
\centering
\renewcommand{\arraystretch}{1.3}
\begin{tabularx}{\textwidth}{|>{\raggedright\arraybackslash\bfseries}p{4.5cm}|>{\raggedright\arraybackslash}X|}
\hline
\rowcolor{tableheader}
\multicolumn{2}{|c|}{\textbf{Đặc tả Use Case: Tạo Page}} \\ \hline
\textbf{Mã Use Case} & UC-14 \\ \hline
\textbf{Tên Use Case} & Tạo Page \\ \hline
\textbf{Actor chính} & Mangaka \\ \hline
\textbf{Mô tả} & Cho phép Mangaka tạo các trang trống (Page) trong Chapter để chuẩn bị phân công công việc cho Assistant. \\ \hline
\textbf{Tiền điều kiện (Pre-conditions)} & Mangaka đăng nhập sở hữu bộ truyện chứa Chapter tương ứng. \\ \hline
\textbf{Luồng sự kiện chính (Basic Flow)} & 
\begin{minipage}[t]{\linewidth}
    \begin{enumerate}[leftmargin=*]
        \item Mangaka vào giao diện chi tiết Chapter.
        \item Mangaka chọn 'Thêm Page'.
        \item Hệ thống tự động sinh số thứ tự trang tiếp theo hoặc cho phép nhập số trang hàng loạt (Batch create).
        \item Mangaka nhấn 'Lưu'.
        \item Hệ thống tạo các bản ghi Page tương ứng trong cơ sở dữ liệu.
    \end{enumerate}
\end{minipage} \\ \hline
\textbf{Luồng rẽ nhánh/Ngoại lệ (Alternate/Exception)} & 
\begin{minipage}[t]{\linewidth}
    \begin{itemize}[leftmargin=*]
        \item \textbf{A1}: Số lượng Page tạo ra vượt mức 100 trang/chapter. Hệ thống cảnh báo quá tải và yêu cầu xác nhận thêm.
    \end{itemize}
\end{minipage} \\ \hline
\textbf{Hậu điều kiện (Post-conditions)} & Các trang mới được tạo với trạng thái 'Trống' (Blank), sẵn sàng cho phân công. \\ \hline
\end{tabularx}
\caption{Đặc tả UC-14: Tạo Page}
\label{tab:uc14}
\end{table}

\subsubsection{UC-15: Cập nhật trạng thái Page}\label{subsubsec:uc15_update_page}
\begin{table}[H]
\centering
\renewcommand{\arraystretch}{1.3}
\begin{tabularx}{\textwidth}{|>{\raggedright\arraybackslash\bfseries}p{4.5cm}|>{\raggedright\arraybackslash}X|}
\hline
\rowcolor{tableheader}
\multicolumn{2}{|c|}{\textbf{Đặc tả Use Case: Cập nhật trạng thái Page}} \\ \hline
\textbf{Mã Use Case} & UC-15 \\ \hline
\textbf{Tên Use Case} & Cập nhật trạng thái Page \\ \hline
\textbf{Actor chính} & Mangaka, Assistant \\ \hline
\textbf{Mô tả} & Cập nhật trạng thái của trang trong quy trình vẽ (Draft, Inked, Toned, Completed). \\ \hline
\textbf{Tiền điều kiện (Pre-conditions)} & Page đã tồn tại và người dùng có quyền chỉnh sửa Page này (được phân công hoặc là Mangaka). \\ \hline
\textbf{Luồng sự kiện chính (Basic Flow)} & 
\begin{minipage}[t]{\linewidth}
    \begin{enumerate}[leftmargin=*]
        \item Người dùng chọn Page cụ thể.
        \item Chọn thay đổi trạng thái từ menu thả xuống (Dropdown).
        \item Hệ thống cập nhật trạng thái mới của Page.
        \item Hệ thống tự động tính toán lại tiến độ tổng của Chapter và lưu lại lịch sử thay đổi.
    \end{enumerate}
\end{minipage} \\ \hline
\textbf{Luồng rẽ nhánh/Ngoại lệ (Alternate/Exception)} & 
\begin{minipage}[t]{\linewidth}
    \begin{itemize}[leftmargin=*]
        \item \textbf{A1}: Người dùng không có quyền (không được phân công). Hệ thống chặn thao tác và hiển thị 'Bạn không có quyền chuyển trạng thái trang này'.
    \end{itemize}
\end{minipage} \\ \hline
\textbf{Hậu điều kiện (Post-conditions)} & Trạng thái Page được lưu vết (Audit log) và tiến độ Chapter thay đổi tương ứng. \\ \hline
\end{tabularx}
\caption{Đặc tả UC-15: Cập nhật trạng thái Page}
\label{tab:uc15}
\end{table}

\subsection{Quản lý Task (Task Assignment \& Tracking)}\label{subsec:srs_task_mgmt}
\subsubsection{UC-16: Tạo và phân công Task}\label{subsubsec:uc16_create_task}
\begin{table}[H]
\centering
\renewcommand{\arraystretch}{1.3}
\begin{tabularx}{\textwidth}{|>{\raggedright\arraybackslash\bfseries}p{4.5cm}|>{\raggedright\arraybackslash}X|}
\hline
\rowcolor{tableheader}
\multicolumn{2}{|c|}{\textbf{Đặc tả Use Case: Tạo và phân công Task}} \\ \hline
\textbf{Mã Use Case} & UC-16 \\ \hline
\textbf{Tên Use Case} & Tạo và phân công Task \\ \hline
\textbf{Actor chính} & Mangaka \\ \hline
\textbf{Mô tả} & Tạo công việc cụ thể (vẽ background, đổ bóng, tone...) và giao cho Assistant. \\ \hline
\textbf{Tiền điều kiện (Pre-conditions)} & Mangaka là người quản lý dự án, danh sách Assistant đã được thêm vào nhóm làm việc. \\ \hline
\textbf{Luồng sự kiện chính (Basic Flow)} & 
\begin{minipage}[t]{\linewidth}
    \begin{enumerate}[leftmargin=*]
        \item Mangaka chọn Page hoặc Chapter.
        \item Chọn 'Tạo Task mới'.
        \item Nhập mô tả công việc (giới hạn 500 ký tự), chọn loại công việc, và ấn định hạn chót (Deadline).
        \item Chọn Assistant từ danh sách thành viên nhóm.
        \item Nhấn 'Phân công'.
        \item Hệ thống lưu Task, kích hoạt gửi Email/Push Notification cho Assistant.
    \end{enumerate}
\end{minipage} \\ \hline
\textbf{Luồng rẽ nhánh/Ngoại lệ (Alternate/Exception)} & 
\begin{minipage}[t]{\linewidth}
    \begin{itemize}[leftmargin=*]
        \item \textbf{A1}: Assistant đang tồn đọng quá 5 Task chưa hoàn thành. Hệ thống hiện cảnh báo 'Assistant đang quá tải' nhưng vẫn cho phép phân công nếu Mangaka xác nhận.
    \end{itemize}
\end{minipage} \\ \hline
\textbf{Hậu điều kiện (Post-conditions)} & Task chuyển sang trạng thái 'Assigned', Assistant nhận được thông báo ngay lập tức. \\ \hline
\end{tabularx}
\caption{Đặc tả UC-16: Tạo và phân công Task}
\label{tab:uc16}
\end{table}

\subsubsection{UC-17: Cập nhật trạng thái Task}\label{subsubsec:uc17_update_task}
\begin{table}[H]
\centering
\renewcommand{\arraystretch}{1.3}
\begin{tabularx}{\textwidth}{|>{\raggedright\arraybackslash\bfseries}p{4.5cm}|>{\raggedright\arraybackslash}X|}
\hline
\rowcolor{tableheader}
\multicolumn{2}{|c|}{\textbf{Đặc tả Use Case: Cập nhật trạng thái Task}} \\ \hline
\textbf{Mã Use Case} & UC-17 \\ \hline
\textbf{Tên Use Case} & Cập nhật trạng thái Task \\ \hline
\textbf{Actor chính} & Assistant \\ \hline
\textbf{Mô tả} & Assistant cập nhật tình trạng công việc mình đang làm trên hệ thống Kanban. \\ \hline
\textbf{Tiền điều kiện (Pre-conditions)} & Assistant đã đăng nhập và là người được chỉ định thực hiện Task đó. \\ \hline
\textbf{Luồng sự kiện chính (Basic Flow)} & 
\begin{minipage}[t]{\linewidth}
    \begin{enumerate}[leftmargin=*]
        \item Assistant vào danh sách 'Task của tôi'.
        \item Kéo thả Task (hoặc bấm chọn) để thay đổi trạng thái sang 'In Progress', 'Ready for Review' hoặc 'Completed'.
        \item Hệ thống ghi nhận trạng thái, thời gian cập nhật.
        \item Thông báo tự động gửi đến Mangaka nếu trạng thái là 'Ready for Review'.
    \end{enumerate}
\end{minipage} \\ \hline
\textbf{Luồng rẽ nhánh/Ngoại lệ (Alternate/Exception)} & 
\begin{minipage}[t]{\linewidth}
    \begin{itemize}[leftmargin=*]
        \item \textbf{A1}: Task thuộc về một Series đang bị tạm ngưng (Hiatus). Hệ thống khóa Task và thông báo 'Dự án đang đóng băng, không thể cập nhật'.
    \end{itemize}
\end{minipage} \\ \hline
\textbf{Hậu điều kiện (Post-conditions)} & Trạng thái Task được thay đổi và lịch sử thao tác được lưu trữ. \\ \hline
\end{tabularx}
\caption{Đặc tả UC-17: Cập nhật trạng thái Task}
\label{tab:uc17}
\end{table}

\subsubsection{UC-18: Theo dõi tiến độ Task}\label{subsubsec:uc18_track_task}
\begin{table}[H]
\centering
\renewcommand{\arraystretch}{1.3}
\begin{tabularx}{\textwidth}{|>{\raggedright\arraybackslash\bfseries}p{4.5cm}|>{\raggedright\arraybackslash}X|}
\hline
\rowcolor{tableheader}
\multicolumn{2}{|c|}{\textbf{Đặc tả Use Case: Theo dõi tiến độ Task}} \\ \hline
\textbf{Mã Use Case} & UC-18 \\ \hline
\textbf{Tên Use Case} & Theo dõi tiến độ Task \\ \hline
\textbf{Actor chính} & Mangaka \\ \hline
\textbf{Mô tả} & Mangaka xem danh sách toàn bộ Task để đánh giá hiệu suất của nhóm. \\ \hline
\textbf{Tiền điều kiện (Pre-conditions)} & Mangaka có quyền quản lý nhóm làm việc (Team management). \\ \hline
\textbf{Luồng sự kiện chính (Basic Flow)} & 
\begin{minipage}[t]{\linewidth}
    \begin{enumerate}[leftmargin=*]
        \item Mangaka mở Dashboard của Series.
        \item Xem danh sách Task dưới dạng Kanban Board (To Do, In Progress, Review, Done).
        \item Sử dụng bộ lọc (Filter) theo Assistant, loại công việc, hoặc theo Deadline gần nhất.
        \item Hệ thống truy vấn và hiển thị các Task đáp ứng tiêu chí lọc ngay lập tức.
    \end{enumerate}
\end{minipage} \\ \hline
\textbf{Luồng rẽ nhánh/Ngoại lệ (Alternate/Exception)} & 
\begin{minipage}[t]{\linewidth}
    \begin{itemize}[leftmargin=*]
        \item \textbf{A1}: Chưa có Task nào được tạo. Hệ thống hiện màn hình trống kèm giao diện hướng dẫn (Onboarding) và nút 'Tạo Task đầu tiên'.
    \end{itemize}
\end{minipage} \\ \hline
\textbf{Hậu điều kiện (Post-conditions)} & Mangaka nắm bắt được nút thắt cổ chai (bottleneck) của nhóm. \\ \hline
\end{tabularx}
\caption{Đặc tả UC-18: Theo dõi tiến độ Task}
\label{tab:uc18}
\end{table}

\subsubsection{UC-19: Hủy Task}\label{subsubsec:uc19_cancel_task}
\begin{table}[H]
\centering
\renewcommand{\arraystretch}{1.3}
\begin{tabularx}{\textwidth}{|>{\raggedright\arraybackslash\bfseries}p{4.5cm}|>{\raggedright\arraybackslash}X|}
\hline
\rowcolor{tableheader}
\multicolumn{2}{|c|}{\textbf{Đặc tả Use Case: Hủy Task}} \\ \hline
\textbf{Mã Use Case} & UC-19 \\ \hline
\textbf{Tên Use Case} & Hủy Task \\ \hline
\textbf{Actor chính} & Mangaka \\ \hline
\textbf{Mô tả} & Hủy bỏ một công việc đã phân công do thay đổi kế hoạch kịch bản. \\ \hline
\textbf{Tiền điều kiện (Pre-conditions)} & Task chưa chuyển sang trạng thái 'Completed' hoặc 'Reviewed'. \\ \hline
\textbf{Luồng sự kiện chính (Basic Flow)} & 
\begin{minipage}[t]{\linewidth}
    \begin{enumerate}[leftmargin=*]
        \item Mangaka chọn Task cần hủy.
        \item Chọn 'Hủy Task', hệ thống yêu cầu nhập lý do bắt buộc.
        \item Mangaka nhập lý do và xác nhận.
        \item Hệ thống đánh dấu Task là 'Cancelled' và gửi thông báo cho Assistant.
    \end{enumerate}
\end{minipage} \\ \hline
\textbf{Luồng rẽ nhánh/Ngoại lệ (Alternate/Exception)} & 
\begin{minipage}[t]{\linewidth}
    \begin{itemize}[leftmargin=*]
        \item \textbf{A1}: Task đã ở trạng thái Completed. Nút Hủy bị mờ (disabled), hệ thống gợi ý 'Tạo Task mới để sửa đổi thay vì hủy'.
    \end{itemize}
\end{minipage} \\ \hline
\textbf{Hậu điều kiện (Post-conditions)} & Task bị vô hiệu hóa, giảm khối lượng công việc hiện tại của Assistant. \\ \hline
\end{tabularx}
\caption{Đặc tả UC-19: Hủy Task}
\label{tab:uc19}
\end{table}

\subsection{Quản lý Submission (Submission Management)}\label{subsec:srs_submission_mgmt}
\subsubsection{UC-20: Nộp Submission}\label{subsubsec:uc20_submit}
\begin{table}[H]
\centering
\renewcommand{\arraystretch}{1.3}
\begin{tabularx}{\textwidth}{|>{\raggedright\arraybackslash\bfseries}p{4.5cm}|>{\raggedright\arraybackslash}X|}
\hline
\rowcolor{tableheader}
\multicolumn{2}{|c|}{\textbf{Đặc tả Use Case: Nộp Submission}} \\ \hline
\textbf{Mã Use Case} & UC-20 \\ \hline
\textbf{Tên Use Case} & Nộp Submission \\ \hline
\textbf{Actor chính} & Assistant \\ \hline
\textbf{Mô tả} & Nộp kết quả công việc (file ảnh/tài liệu) lên Cloud để Mangaka đánh giá. \\ \hline
\textbf{Tiền điều kiện (Pre-conditions)} & Task đang ở trạng thái 'In Progress' và thuộc quyền xử lý của Assistant. \\ \hline
\textbf{Luồng sự kiện chính (Basic Flow)} & 
\begin{minipage}[t]{\linewidth}
    \begin{enumerate}[leftmargin=*]
        \item Assistant chọn Task tương ứng.
        \item Chọn chức năng 'Nộp kết quả (Submission)'.
        \item Tải lên các file kết quả (Chỉ chấp nhận JPEG, PNG, PSD; dung lượng tối đa 20MB/file).
        \item Viết ghi chú đính kèm.
        \item Nhấn 'Gửi'. Hệ thống lưu file lên Cloud Storage, chuyển trạng thái Task sang 'Reviewing' và gửi Notification cho Mangaka.
    \end{enumerate}
\end{minipage} \\ \hline
\textbf{Luồng rẽ nhánh/Ngoại lệ (Alternate/Exception)} & 
\begin{minipage}[t]{\linewidth}
    \begin{itemize}[leftmargin=*]
        \item \textbf{A1}: File sai định dạng hoặc quá dung lượng 20MB. Hệ thống từ chối tải lên, bôi đỏ lỗi và hướng dẫn Assistant nén file hoặc đổi định dạng.
    \end{itemize}
\end{minipage} \\ \hline
\textbf{Hậu điều kiện (Post-conditions)} & Submission được ghi nhận an toàn, phiên bản (version) được cập nhật. \\ \hline
\end{tabularx}
\caption{Đặc tả UC-20: Nộp Submission}
\label{tab:uc20}
\end{table}

\subsubsection{UC-21: Xem chi tiết Submission}\label{subsubsec:uc21_view_submission}
\begin{table}[H]
\centering
\renewcommand{\arraystretch}{1.3}
\begin{tabularx}{\textwidth}{|>{\raggedright\arraybackslash\bfseries}p{4.5cm}|>{\raggedright\arraybackslash}X|}
\hline
\rowcolor{tableheader}
\multicolumn{2}{|c|}{\textbf{Đặc tả Use Case: Xem chi tiết Submission}} \\ \hline
\textbf{Mã Use Case} & UC-21 \\ \hline
\textbf{Tên Use Case} & Xem chi tiết Submission \\ \hline
\textbf{Actor chính} & Mangaka, Tantou Editor \\ \hline
\textbf{Mô tả} & Xem bản thảo chất lượng cao đã nộp để đưa ra nhận xét trực quan. \\ \hline
\textbf{Tiền điều kiện (Pre-conditions)} & Tồn tại ít nhất 1 Submission và tài khoản có quyền truy cập dữ liệu của Task. \\ \hline
\textbf{Luồng sự kiện chính (Basic Flow)} & 
\begin{minipage}[t]{\linewidth}
    \begin{enumerate}[leftmargin=*]
        \item Người dùng nhấn vào link trong thông báo hoặc mở chi tiết Task.
        \item Hệ thống render file bản thảo (nếu là PSD sẽ render bản preview) tích hợp công cụ thu phóng ảnh.
        \item Người dùng xem nội dung hình ảnh và đọc các ghi chú đi kèm.
    \end{enumerate}
\end{minipage} \\ \hline
\textbf{Luồng rẽ nhánh/Ngoại lệ (Alternate/Exception)} & 
\begin{minipage}[t]{\linewidth}
    \begin{itemize}[leftmargin=*]
        \item \textbf{A1}: File bị lỗi corrupt (hỏng) không thể render. Hệ thống thông báo lỗi hiển thị và cung cấp nút 'Tải file gốc về máy' để xem cục bộ.
    \end{itemize}
\end{minipage} \\ \hline
\textbf{Hậu điều kiện (Post-conditions)} & Người dùng tiếp nhận đầy đủ kết quả công việc của Assistant. \\ \hline
\end{tabularx}
\caption{Đặc tả UC-21: Xem chi tiết Submission}
\label{tab:uc21}
\end{table}

\subsubsection{UC-22: Xem lịch sử Submission}\label{subsubsec:uc22_submission_history}
\begin{table}[H]
\centering
\renewcommand{\arraystretch}{1.3}
\begin{tabularx}{\textwidth}{|>{\raggedright\arraybackslash\bfseries}p{4.5cm}|>{\raggedright\arraybackslash}X|}
\hline
\rowcolor{tableheader}
\multicolumn{2}{|c|}{\textbf{Đặc tả Use Case: Xem lịch sử Submission}} \\ \hline
\textbf{Mã Use Case} & UC-22 \\ \hline
\textbf{Tên Use Case} & Xem lịch sử Submission \\ \hline
\textbf{Actor chính} & Tất cả \\ \hline
\textbf{Mô tả} & Truy xuất và xem lại toàn bộ các lần nộp bản thảo (phiên bản) trước đó của một Task/Page. \\ \hline
\textbf{Tiền điều kiện (Pre-conditions)} & Task/Page có từ 1 lần Submission trở lên. \\ \hline
\textbf{Luồng sự kiện chính (Basic Flow)} & 
\begin{minipage}[t]{\linewidth}
    \begin{enumerate}[leftmargin=*]
        \item Người dùng chọn 'Lịch sử nộp bài' (Version History).
        \item Hệ thống truy xuất Database, liệt kê danh sách các version theo dòng thời gian (Timeline).
        \item Người dùng có thể nhấn vào từng version để xem bản lưu trước đó.
    \end{enumerate}
\end{minipage} \\ \hline
\textbf{Luồng rẽ nhánh/Ngoại lệ (Alternate/Exception)} & 
\begin{minipage}[t]{\linewidth}
    \begin{itemize}[leftmargin=*]
        \item \textbf{A1}: Phiên bản cũ đã bị xóa khỏi lưu trữ để tiết kiệm dung lượng. Hệ thống hiển thị 'Phiên bản này đã lưu trữ ngoại tuyến (Archived)'.
    \end{itemize}
\end{minipage} \\ \hline
\textbf{Hậu điều kiện (Post-conditions)} & Tính minh bạch của quá trình chỉnh sửa bản thảo được đảm bảo. \\ \hline
\end{tabularx}
\caption{Đặc tả UC-22: Xem lịch sử Submission}
\label{tab:uc22}
\end{table}

\subsection{Quản lý Review (Review Management)}\label{subsec:srs_review_mgmt}
\subsubsection{UC-23: Tạo Review cho Submission}\label{subsubsec:uc23_create_review}
\begin{table}[H]
\centering
\renewcommand{\arraystretch}{1.3}
\begin{tabularx}{\textwidth}{|>{\raggedright\arraybackslash\bfseries}p{4.5cm}|>{\raggedright\arraybackslash}X|}
\hline
\rowcolor{tableheader}
\multicolumn{2}{|c|}{\textbf{Đặc tả Use Case: Tạo Review cho Submission}} \\ \hline
\textbf{Mã Use Case} & UC-23 \\ \hline
\textbf{Tên Use Case} & Tạo Review cho Submission \\ \hline
\textbf{Actor chính} & Tantou Editor, Mangaka \\ \hline
\textbf{Mô tả} & Thêm nhận xét, vẽ khoanh vùng lỗi trực tiếp lên ảnh bản thảo và yêu cầu chỉnh sửa. \\ \hline
\textbf{Tiền điều kiện (Pre-conditions)} & Đang xem giao diện chi tiết một Submission chưa được Approve. \\ \hline
\textbf{Luồng sự kiện chính (Basic Flow)} & 
\begin{minipage}[t]{\linewidth}
    \begin{enumerate}[leftmargin=*]
        \item Người dùng chọn công cụ vẽ (highlight, khoanh vùng, bút đỏ) trên trình duyệt.
        \item Vẽ đánh dấu trực tiếp lên các điểm cần sửa trên file ảnh.
        \item Nhập comment nhận xét (văn bản) cho từng vùng đánh dấu.
        \item Chọn 'Approve' (Chấp nhận) hoặc 'Request Changes' (Yêu cầu sửa).
        \item Hệ thống lưu tọa độ đánh dấu, nội dung text và gửi email/push notification cho tác giả/trợ lý.
    \end{enumerate}
\end{minipage} \\ \hline
\textbf{Luồng rẽ nhánh/Ngoại lệ (Alternate/Exception)} & 
\begin{minipage}[t]{\linewidth}
    \begin{itemize}[leftmargin=*]
        \item \textbf{A1}: Mất kết nối mạng khi đang khoanh vùng lỗi. Hệ thống tự động lưu nháp (Auto-save) vào LocalStorage của trình duyệt và cảnh báo 'Đang ngoại tuyến'.
    \end{itemize}
\end{minipage} \\ \hline
\textbf{Hậu điều kiện (Post-conditions)} & Tọa độ Annotation (được lưu trữ tại cột \texttt{annotations} trong cơ sở dữ liệu) và kết quả Review được lưu trữ vĩnh viễn, trạng thái Task thay đổi tương ứng. \\ \hline
\end{tabularx}
\caption{Đặc tả UC-23: Tạo Review cho Submission}
\label{tab:uc23}
\end{table}

\subsubsection{UC-24: Xem chi tiết Review}\label{subsubsec:uc24_view_review}
\begin{table}[H]
\centering
\renewcommand{\arraystretch}{1.3}
\begin{tabularx}{\textwidth}{|>{\raggedright\arraybackslash\bfseries}p{4.5cm}|>{\raggedright\arraybackslash}X|}
\hline
\rowcolor{tableheader}
\multicolumn{2}{|c|}{\textbf{Đặc tả Use Case: Xem chi tiết Review}} \\ \hline
\textbf{Mã Use Case} & UC-24 \\ \hline
\textbf{Tên Use Case} & Xem chi tiết Review \\ \hline
\textbf{Actor chính} & Assistant, Mangaka \\ \hline
\textbf{Mô tả} & Người thực hiện công việc xem lại các điểm cần sửa chữa do Editor/Mangaka đánh dấu. \\ \hline
\textbf{Tiền điều kiện (Pre-conditions)} & Submission đã được Review và có ít nhất 1 nhận xét hoặc đánh dấu. \\ \hline
\textbf{Luồng sự kiện chính (Basic Flow)} & 
\begin{minipage}[t]{\linewidth}
    \begin{enumerate}[leftmargin=*]
        \item Người dùng mở Submission đã bị gắn cờ 'Request Changes'.
        \item Hệ thống tải tọa độ Annotation, hiển thị ảnh gốc xếp chồng (overlay) các vùng khoanh đỏ của reviewer.
        \item Người dùng trỏ chuột (hover) vào vùng khoanh đỏ để đọc bình luận chi tiết tương ứng.
    \end{enumerate}
\end{minipage} \\ \hline
\textbf{Luồng rẽ nhánh/Ngoại lệ (Alternate/Exception)} & 
\begin{minipage}[t]{\linewidth}
    \begin{itemize}[leftmargin=*]
        \item \textbf{A1}: Trình duyệt không hỗ trợ Canvas render. Hệ thống tự động chuyển sang chế độ hiển thị danh sách nhận xét dạng text thuần túy.
    \end{itemize}
\end{minipage} \\ \hline
\textbf{Hậu điều kiện (Post-conditions)} & Người thực hiện định vị chính xác vị trí lỗi cần khắc phục trên bức tranh. \\ \hline
\end{tabularx}
\caption{Đặc tả UC-24: Xem chi tiết Review}
\label{tab:uc24}
\end{table}

\subsubsection{UC-25: Phản hồi yêu cầu chỉnh sửa}\label{subsubsec:uc25_respond_review}
\begin{table}[H]
\centering
\renewcommand{\arraystretch}{1.3}
\begin{tabularx}{\textwidth}{|>{\raggedright\arraybackslash\bfseries}p{4.5cm}|>{\raggedright\arraybackslash}X|}
\hline
\rowcolor{tableheader}
\multicolumn{2}{|c|}{\textbf{Đặc tả Use Case: Phản hồi yêu cầu chỉnh sửa}} \\ \hline
\textbf{Mã Use Case} & UC-25 \\ \hline
\textbf{Tên Use Case} & Phản hồi yêu cầu chỉnh sửa \\ \hline
\textbf{Actor chính} & Assistant, Mangaka \\ \hline
\textbf{Mô tả} & Gửi lại file mới (Revision) sau khi đã sửa lỗi theo yêu cầu của Review. \\ \hline
\textbf{Tiền điều kiện (Pre-conditions)} & Tồn tại một Review ở trạng thái 'Request Changes' yêu cầu người dùng phải hành động. \\ \hline
\textbf{Luồng sự kiện chính (Basic Flow)} & 
\begin{minipage}[t]{\linewidth}
    \begin{enumerate}[leftmargin=*]
        \item Assistant hoàn tất sửa lỗi trên phần mềm vẽ cục bộ.
        \item Vào mục Review, chọn chức năng 'Nộp bản sửa' (Submit Revision).
        \item Tải file bản vá (patch) lên và tích vào các checkbox đánh dấu 'Đã sửa' cho từng comment của Reviewer.
        \item Hệ thống tăng version của Submission, liên kết bản mới này với Review cũ và thông báo cho Reviewer kiểm tra lại.
    \end{enumerate}
\end{minipage} \\ \hline
\textbf{Luồng rẽ nhánh/Ngoại lệ (Alternate/Exception)} & 
\begin{minipage}[t]{\linewidth}
    \begin{itemize}[leftmargin=*]
        \item \textbf{A1}: Thuật toán Hash phát hiện file gửi lên giống hệt file cũ. Hệ thống từ chối tải lên và yêu cầu 'Vui lòng nộp file đã có chỉnh sửa'.
    \end{itemize}
\end{minipage} \\ \hline
\textbf{Hậu điều kiện (Post-conditions)} & Chu trình duyệt bài lặp lại, đảm bảo kiểm soát chất lượng. \\ \hline
\end{tabularx}
\caption{Đặc tả UC-25: Phản hồi yêu cầu chỉnh sửa}
\label{tab:uc25}
\end{table}

\subsection{Xét duyệt \& Xuất bản (Approval \& Publishing)}\label{subsec:srs_publishing_mgmt}
\subsubsection{UC-26: Xét duyệt Series mới}\label{subsubsec:uc26_approve_series}
\begin{table}[H]
\centering
\renewcommand{\arraystretch}{1.3}
\begin{tabularx}{\textwidth}{|>{\raggedright\arraybackslash\bfseries}p{4.5cm}|>{\raggedright\arraybackslash}X|}
\hline
\rowcolor{tableheader}
\multicolumn{2}{|c|}{\textbf{Đặc tả Use Case: Xét duyệt Series mới}} \\ \hline
\textbf{Mã Use Case} & UC-26 \\ \hline
\textbf{Tên Use Case} & Xét duyệt Series mới \\ \hline
\textbf{Actor chính} & Editorial Board \\ \hline
\textbf{Mô tả} & Từng thành viên hội đồng tiến hành xem xét kịch bản, bỏ phiếu và phê duyệt một đề xuất Series mới. \\ \hline
\textbf{Tiền điều kiện (Pre-conditions)} & Tài khoản thuộc nhóm Editorial Board, tồn tại Series ở trạng thái 'planning' và đã được nộp đề xuất (publish\_type = 'submitted'). \\ \hline
\textbf{Luồng sự kiện chính (Basic Flow)} & 
\begin{minipage}[t]{\linewidth}
    \begin{enumerate}[leftmargin=*]
        \item Hội đồng xem xét tài liệu đề xuất và file kịch bản (proposal\_file) của Series.
        \item Từng thành viên thực hiện bỏ phiếu (Đồng ý - Approve hoặc Từ chối - Reject) trực tiếp trên giao diện hệ thống.
        \item Hệ thống ghi nhận phiếu bầu, gửi thông báo cập nhật tiến trình cho Mangaka và các thành viên khác trong Hội đồng.
        \item Khi nhận đầy đủ phiếu bầu từ tất cả thành viên (100\%), hệ thống gửi thông báo cho các thành viên Board sẵn sàng tiến hành chốt quyết định phê duyệt.
    \end{enumerate}
\end{minipage} \\ \hline
\textbf{Luồng rẽ nhánh/Ngoại lệ (Alternate/Exception)} & 
\begin{minipage}[t]{\linewidth}
    \begin{itemize}[leftmargin=*]
        \item \textbf{A1}: Chưa nhận đầy đủ phiếu bầu từ các thành viên Hội đồng. Hệ thống sẽ khóa chức năng chốt quyết định phê duyệt trạng thái Series.
        \item \textbf{A2}: Tỷ lệ tán thành phiếu bầu của Hội đồng chưa đạt tối thiểu 50\%. Hệ thống sẽ chặn không cho phép phê duyệt sang trạng thái Đang triển khai (Ongoing).
    \end{itemize}
\end{minipage} \\ \hline
\textbf{Hậu điều kiện (Post-conditions)} & Ghi nhận đầy đủ phiếu biểu quyết, thông báo kết quả tiến trình bỏ phiếu của Hội đồng. \\ \hline
\end{tabularx}
\caption{Đặc tả UC-26: Xét duyệt Series mới}
\label{tab:uc26}
\end{table}

\subsubsection{UC-27: Phê duyệt Chapter}\label{subsubsec:uc27_publish_chapter}
\begin{table}[H]
\centering
\renewcommand{\arraystretch}{1.3}
\begin{tabularx}{\textwidth}{|>{\raggedright\arraybackslash\bfseries}p{4.5cm}|>{\raggedright\arraybackslash}X|}
\hline
\rowcolor{tableheader}
\multicolumn{2}{|c|}{\textbf{Đặc tả Use Case: Phê duyệt Chapter}} \\ \hline
\textbf{Mã Use Case} & UC-27 \\ \hline
\textbf{Tên Use Case} & Phê duyệt Chapter \\ \hline
\textbf{Actor chính} & Tantou Editor \\ \hline
\textbf{Mô tả} & Tantou Editor phê duyệt chương truyện đã hoàn thành sau khi kiểm duyệt bản thảo cuối cùng đạt chất lượng. \\ \hline
\textbf{Tiền điều kiện (Pre-conditions)} & Chapter hoàn thành các Task vẽ trang truyện, Mangaka nộp bản thảo Chapter hoàn chỉnh lên hệ thống (Chương ở trạng thái 'reviewing\_final'). \\ \hline
\textbf{Luồng sự kiện chính (Basic Flow)} & 
\begin{minipage}[t]{\linewidth}
    \begin{enumerate}[leftmargin=*]
        \item Tantou Editor truy cập vào trang kiểm duyệt bản thảo chương truyện.
        \item Kiểm tra chất lượng hình ảnh vẽ và nội dung của từng trang truyện.
        \item Chọn kết quả phê duyệt là "Approved".
        \item Hệ thống tự động chuyển trạng thái của Chapter sang 'approved', đồng thời khóa cứng toàn bộ các trang truyện và các Task liên quan để chuẩn bị xuất bản.
        \item Hệ thống gửi thông báo phê duyệt thành công cho Mangaka.
    \end{enumerate}
\end{minipage} \\ \hline
\textbf{Luồng rẽ nhánh/Ngoại lệ (Alternate/Exception)} & 
\begin{minipage}[t]{\linewidth}
    \begin{itemize}[leftmargin=*]
        \item \textbf{A1}: Bản thảo chương chưa đạt yêu cầu. Editor chọn "Từ chối" (Rejected), ghi chú các lỗi cần sửa và trả lại. Hệ thống tự động chuyển trạng thái Chapter về 'drawing' để Mangaka và trợ lý chỉnh sửa lại.
    \end{itemize}
\end{minipage} \\ \hline
\textbf{Hậu điều kiện (Post-conditions)} & Chapter được chuyển trạng thái thành 'approved' và khóa chỉnh sửa, sẵn sàng phát hành. \\ \hline
\end{tabularx}
\caption{Đặc tả UC-27: Phê duyệt Chapter}
\label{tab:uc27}
\end{table}

\subsubsection{UC-28: Cập nhật trạng thái Series}\label{subsubsec:uc28_cease_publish}
\begin{table}[H]
\centering
\renewcommand{\arraystretch}{1.3}
\begin{tabularx}{\textwidth}{|>{\raggedright\arraybackslash\bfseries}p{4.5cm}|>{\raggedright\arraybackslash}X|}
\hline
\rowcolor{tableheader}
\multicolumn{2}{|c|}{\textbf{Đặc tả Use Case: Cập nhật trạng thái Series}} \\ \hline
\textbf{Mã Use Case} & UC-28 \\ \hline
\textbf{Tên Use Case} & Cập nhật trạng thái Series \\ \hline
\textbf{Actor chính} & Editorial Board \\ \hline
\textbf{Mô tả} & Cập nhật trạng thái hoạt động (Ongoing, Completed, Canceled, Suspended) và lịch xuất bản của bộ truyện. \\ \hline
\textbf{Tiền điều kiện (Pre-conditions)} & Tài khoản thuộc nhóm Editorial Board. \\ \hline
\textbf{Luồng sự kiện chính (Basic Flow)} & 
\begin{minipage}[t]{\linewidth}
    \begin{enumerate}[leftmargin=*]
        \item Hội đồng truy cập vào trang Quản lý xuất bản bộ truyện.
        \item Chọn bộ truyện cần cập nhật trạng thái.
        \item Chọn trạng thái mới và lịch xuất bản (Hàng tuần/Hàng tháng) phù hợp.
        \item Hệ thống kiểm tra các ràng buộc nghiệp vụ (Business Rules).
        \item Nếu hợp lệ, hệ thống cập nhật trạng thái mới, ghi nhật ký hoạt động (\texttt{SystemLog}) và gửi thông báo kết quả cho Mangaka.
    \end{enumerate}
\end{minipage} \\ \hline
\textbf{Luồng rẽ nhánh/Ngoại lệ (Alternate/Exception)} & 
\begin{minipage}[t]{\linewidth}
    \begin{itemize}[leftmargin=*]
        \item \textbf{A1}: Chọn trạng thái Hoàn thành (Completed) nhưng vẫn còn chapter chưa hoàn thành hoặc chưa có chapter nào là chương cuối (\texttt{is\_final = 1}) được duyệt $\rightarrow$ Hệ thống báo lỗi và từ chối cập nhật.
        \item \textbf{A2}: Phê duyệt sang Đang triển khai (Ongoing) nhưng chưa đủ 100\% phiếu bầu, tỉ lệ tán thành $< 50\%$, chưa gán Editor chuyên trách, hoặc bộ truyện đang ở dạng nháp/thiếu kịch bản đề xuất $\rightarrow$ Hệ thống báo lỗi và từ chối cập nhật.
        \item \textbf{A3}: Hạ cấp bộ truyện về trạng thái Kế hoạch (Planning) nhưng đã có chapter được tạo $\rightarrow$ Hệ thống báo lỗi và từ chối cập nhật.
    \end{itemize}
\end{minipage} \\ \hline
\textbf{Hậu điều kiện (Post-conditions)} & Trạng thái hoạt động của Series được cập nhật mới trên hệ thống, ghi nhận nhật ký hoạt động. \\ \hline
\end{tabularx}
\caption{Đặc tả UC-28: Cập nhật trạng thái Series}
\label{tab:uc28}
\end{table}

\subsection{Quản lý Xếp hạng (Ranking Management)}\label{subsec:srs_ranking_mgmt}
\subsubsection{UC-29: Nhập dữ liệu bình chọn}\label{subsubsec:uc29_input_votes}
\begin{table}[H]
\centering
\renewcommand{\arraystretch}{1.3}
\begin{tabularx}{\textwidth}{|>{\raggedright\arraybackslash\bfseries}p{4.5cm}|>{\raggedright\arraybackslash}X|}
\hline
\rowcolor{tableheader}
\multicolumn{2}{|c|}{\textbf{Đặc tả Use Case: Nhập dữ liệu bình chọn}} \\ \hline
\textbf{Mã Use Case} & UC-29 \\ \hline
\textbf{Tên Use Case} & Nhập dữ liệu bình chọn \\ \hline
\textbf{Actor chính} & Editorial Board \\ \hline
\textbf{Mô tả} & Nhập số liệu phiếu bầu của độc giả trong kỳ phát hành để hệ thống tự động tính điểm quy chuẩn và bảng xếp hạng. \\ \hline
\textbf{Tiền điều kiện (Pre-conditions)} & Tài khoản có quyền Editorial Board, có số liệu bình chọn của độc giả sau kỳ phát hành. \\ \hline
\textbf{Luồng sự kiện chính (Basic Flow)} & 
\begin{minipage}[t]{\linewidth}
    \begin{enumerate}[leftmargin=*]
        \item Thành viên Hội đồng vào module Xếp hạng $\rightarrow$ chọn chức năng "Nhập phiếu bầu".
        \item Chọn Ngày bắt đầu kỳ đánh giá (Period Start Date).
        \item Nhập số phiếu bình chọn thực tế của độc giả cho từng Series đang hoạt động.
        \item Nhấn "Tính toán \& Lưu".
        \item Hệ thống tự động xóa dữ liệu xếp hạng cũ cùng kỳ (nếu có), tính toán Điểm số quy chuẩn (Score) từ mốc phiếu cao nhất (100 điểm), và xếp hạng (Rank Position) theo chuẩn thi đấu (đồng hạng nếu cùng số phiếu).
        \item Hệ thống lưu thông tin, ghi log hoạt động của Board vào \texttt{SystemLog} và gửi thông báo kết quả cho Mangaka.
    \end{enumerate}
\end{minipage} \\ \hline
\textbf{Luồng rẽ nhánh/Ngoại lệ (Alternate/Exception)} & 
\begin{minipage}[t]{\linewidth}
    \begin{itemize}[leftmargin=*]
        \item \textbf{A1}: Điểm quy chuẩn của bộ truyện dưới 50 điểm $\rightarrow$ Hệ thống tự động gửi thêm thông báo cảnh báo (\texttt{series\_warning}) nguy cơ bị ngừng phát hành cho Mangaka phụ trách.
        \item \textbf{A2}: Chọn ngày đánh giá ở tương lai $\rightarrow$ Hệ thống báo lỗi và từ chối lưu dữ liệu.
        \item \textbf{A3}: Nhập số phiếu bầu âm $\rightarrow$ Hệ thống tự động chuyển số phiếu về 0.
    \end{itemize}
\end{minipage} \\ \hline
\textbf{Hậu điều kiện (Post-conditions)} & Bảng xếp hạng của kỳ phát hành được cập nhật và gửi thông báo kết quả cho các Mangaka. \\ \hline
\end{tabularx}
\caption{Đặc tả UC-29: Nhập dữ liệu bình chọn}
\label{tab:uc29}
\end{table}

\subsubsection{UC-30: Xem bảng xếp hạng}\label{subsubsec:uc30_view_ranking}
\begin{table}[H]
\centering
\renewcommand{\arraystretch}{1.3}
\begin{tabularx}{\textwidth}{|>{\raggedright\arraybackslash\bfseries}p{4.5cm}|>{\raggedright\arraybackslash}X|}
\hline
\rowcolor{tableheader}
\multicolumn{2}{|c|}{\textbf{Đặc tả Use Case: Xem bảng xếp hạng}} \\ \hline
\textbf{Mã Use Case} & UC-30 \\ \hline
\textbf{Tên Use Case} & Xem bảng xếp hạng \\ \hline
\textbf{Actor chính} & Tất cả \\ \hline
\textbf{Mô tả} & Xem thứ hạng hiện tại của các Series trên tạp chí để đánh giá mức độ phổ biến. \\ \hline
\textbf{Tiền điều kiện (Pre-conditions)} & Người dùng đã đăng nhập hệ thống nội bộ. \\ \hline
\textbf{Luồng sự kiện chính (Basic Flow)} & 
\begin{minipage}[t]{\linewidth}
    \begin{enumerate}[leftmargin=*]
        \item Người dùng chọn menu 'Bảng xếp hạng'.
        \item Hệ thống truy xuất và hiển thị danh sách Series sắp xếp theo thứ hạng của kỳ mới nhất.
        \item Giao diện hiển thị mũi tên xanh/đỏ biểu thị xu hướng (tăng/giảm hạng) so với kỳ phát hành liền trước.
    \end{enumerate}
\end{minipage} \\ \hline
\textbf{Luồng rẽ nhánh/Ngoại lệ (Alternate/Exception)} & 
\begin{minipage}[t]{\linewidth}
    \begin{itemize}[leftmargin=*]
        \item \textbf{A1}: Chưa có dữ liệu bình chọn của kỳ hiện tại. Hệ thống hiển thị bảng xếp hạng của kỳ gần nhất kèm dòng thông báo 'Dữ liệu tuần này đang được cập nhật'.
    \end{itemize}
\end{minipage} \\ \hline
\textbf{Hậu điều kiện (Post-conditions)} & Mangaka và Editor nắm được tình hình cạnh tranh và áp lực thị hiếu. \\ \hline
\end{tabularx}
\caption{Đặc tả UC-30: Xem bảng xếp hạng}
\label{tab:uc30}
\end{table}

\subsubsection{UC-31: Xem lịch sử xếp hạng}\label{subsubsec:uc31_ranking_history}
\begin{table}[H]
\centering
\renewcommand{\arraystretch}{1.3}
\begin{tabularx}{\textwidth}{|>{\raggedright\arraybackslash\bfseries}p{4.5cm}|>{\raggedright\arraybackslash}X|}
\hline
\rowcolor{tableheader}
\multicolumn{2}{|c|}{\textbf{Đặc tả Use Case: Xem lịch sử xếp hạng}} \\ \hline
\textbf{Mã Use Case} & UC-31 \\ \hline
\textbf{Tên Use Case} & Xem lịch sử xếp hạng \\ \hline
\textbf{Actor chính} & Tất cả \\ \hline
\textbf{Mô tả} & Xem biểu đồ biến động thứ hạng (Ranking Trend) của một Series cụ thể qua thời gian. \\ \hline
\textbf{Tiền điều kiện (Pre-conditions)} & Người dùng đang xem Dashboard chi tiết của một Series. \\ \hline
\textbf{Luồng sự kiện chính (Basic Flow)} & 
\begin{minipage}[t]{\linewidth}
    \begin{enumerate}[leftmargin=*]
        \item Người dùng chuyển sang tab 'Lịch sử xếp hạng' (Ranking History) của Series.
        \item Hệ thống vẽ biểu đồ đường (Line chart) thể hiện thứ hạng qua các tuần/kỳ.
        \item Người dùng trỏ chuột (hover) vào từng giao điểm để xem chính xác số phiếu bầu và thứ hạng.
    \end{enumerate}
\end{minipage} \\ \hline
\textbf{Luồng rẽ nhánh/Ngoại lệ (Alternate/Exception)} & 
\begin{minipage}[t]{\linewidth}
    \begin{itemize}[leftmargin=*]
        \item \textbf{A1}: Series mới phát hành được 1 kỳ. Biểu đồ chỉ hiển thị dạng 1 điểm (Dot) và thông báo 'Chưa đủ dữ liệu vẽ xu hướng'.
    \end{itemize}
\end{minipage} \\ \hline
\textbf{Hậu điều kiện (Post-conditions)} & Cung cấp phân tích trực quan về vòng đời và mức độ yêu thích của tác phẩm. \\ \hline
\end{tabularx}
\caption{Đặc tả UC-31: Xem lịch sử xếp hạng}
\label{tab:uc31}
\end{table}

\subsection{Quản lý Thông báo (Notification Management)}\label{subsec:srs_notification_mgmt}
\subsubsection{UC-32: Xem danh sách thông báo}\label{subsubsec:uc32_view_notifications}
\begin{table}[H]
\centering
\renewcommand{\arraystretch}{1.3}
\begin{tabularx}{\textwidth}{|>{\raggedright\arraybackslash\bfseries}p{4.5cm}|>{\raggedright\arraybackslash}X|}
\hline
\rowcolor{tableheader}
\multicolumn{2}{|c|}{\textbf{Đặc tả Use Case: Xem danh sách thông báo}} \\ \hline
\textbf{Mã Use Case} & UC-32 \\ \hline
\textbf{Tên Use Case} & Xem danh sách thông báo \\ \hline
\textbf{Actor chính} & Tất cả \\ \hline
\textbf{Mô tả} & Xem các thông báo hệ thống gửi đến tài khoản (nhắc nhở deadline, có task mới, review mới). \\ \hline
\textbf{Tiền điều kiện (Pre-conditions)} & Người dùng đã đăng nhập thành công. \\ \hline
\textbf{Luồng sự kiện chính (Basic Flow)} & 
\begin{minipage}[t]{\linewidth}
    \begin{enumerate}[leftmargin=*]
        \item Người dùng nhấn vào biểu tượng 'Chuông thông báo' trên góc phải màn hình.
        \item Hệ thống truy vấn (Real-time Socket) và xổ xuống menu (dropdown) chứa 5-10 thông báo mới nhất.
        \item Người dùng có thể bấm 'Xem tất cả' để điều hướng sang trang quản lý thông báo chi tiết.
    \end{enumerate}
\end{minipage} \\ \hline
\textbf{Luồng rẽ nhánh/Ngoại lệ (Alternate/Exception)} & 
\begin{minipage}[t]{\linewidth}
    \begin{itemize}[leftmargin=*]
        \item \textbf{A1}: Không có thông báo nào trong cơ sở dữ liệu. Menu thả xuống hiển thị ảnh minh họa (empty state) 'Bạn đã xem hết thông báo mới'.
    \end{itemize}
\end{minipage} \\ \hline
\textbf{Hậu điều kiện (Post-conditions)} & Người dùng nắm bắt được các sự kiện quan trọng cần xử lý gấp. \\ \hline
\end{tabularx}
\caption{Đặc tả UC-32: Xem danh sách thông báo}
\label{tab:uc32}
\end{table}

\subsubsection{UC-33: Đánh dấu đã đọc}\label{subsubsec:uc33_read_notification}
\begin{table}[H]
\centering
\renewcommand{\arraystretch}{1.3}
\begin{tabularx}{\textwidth}{|>{\raggedright\arraybackslash\bfseries}p{4.5cm}|>{\raggedright\arraybackslash}X|}
\hline
\rowcolor{tableheader}
\multicolumn{2}{|c|}{\textbf{Đặc tả Use Case: Đánh dấu đã đọc}} \\ \hline
\textbf{Mã Use Case} & UC-33 \\ \hline
\textbf{Tên Use Case} & Đánh dấu đã đọc \\ \hline
\textbf{Actor chính} & Tất cả \\ \hline
\textbf{Mô tả} & Đánh dấu một hoặc toàn bộ thông báo là đã xem để xóa huy hiệu (badge) báo hiệu chưa đọc. \\ \hline
\textbf{Tiền điều kiện (Pre-conditions)} & Có ít nhất 1 thông báo ở trạng thái chưa đọc (unread). \\ \hline
\textbf{Luồng sự kiện chính (Basic Flow)} & 
\begin{minipage}[t]{\linewidth}
    \begin{enumerate}[leftmargin=*]
        \item Người dùng bấm vào một thông báo cụ thể (Hệ thống tự động kích hoạt chuyển trạng thái thành đã đọc).
        \item Hoặc người dùng bấm nút 'Đánh dấu tất cả đã đọc' (Mark all as read).
        \item Hệ thống cập nhật trường trạng thái trong DB.
        \item Biểu tượng số lượng thông báo mới (đỏ) trên chuông giảm đi tương ứng hoặc biến mất.
    \end{enumerate}
\end{minipage} \\ \hline
\textbf{Luồng rẽ nhánh/Ngoại lệ (Alternate/Exception)} & 
\begin{minipage}[t]{\linewidth}
    \begin{itemize}[leftmargin=*]
        \item \textbf{A1}: Quá trình cập nhật DB thất bại do lỗi máy chủ. UI tự động khôi phục lại trạng thái chưa đọc và hiển thị toast 'Có lỗi xảy ra, thử lại sau'.
    \end{itemize}
\end{minipage} \\ \hline
\textbf{Hậu điều kiện (Post-conditions)} & Giải phóng người dùng khỏi các nhắc nhở không còn cần thiết. \\ \hline
\end{tabularx}
\caption{Đặc tả UC-33: Đánh dấu đã đọc}
\label{tab:uc33}
\end{table}

